\chapter{Introduction}
\label{chap:introduction}

\section{Background}

Investor sentiment plays a well-documented role in shaping financial markets.
Beyond traditional indicators (earnings reports, macroeconomic data, and analyst
forecasts), collective mood and perception can drive price volatility, amplify
trading volume, and create temporary mispricings. These mispricings persist
until information is fully absorbed \citep{tetlock2007giving}. With the
proliferation of digital communication, social media and online communities
have become primary channels through which retail investors express and
exchange market opinions. Reddit,
X (formerly Twitter), and financial news outlets together host millions of daily
discussions that reflect how diverse investor communities interpret and respond to
market developments.

The emergence of Natural Language Processing (NLP) has made it possible to
quantify these qualitative opinions at scale. In particular, domain-adaptive
transformer models such as FinBERT \citep{araci2019finbert} have demonstrated
state-of-the-art performance on financial sentiment benchmarks, capturing nuances
such as ``guidance raised'', ``missed expectations'', and ``strong headwinds'' that
earlier lexicon or bag-of-words approaches simply miss. These advances open a
realistic path toward decision-support tools that fuse public opinion signals
with quantitative market data.

What I noticed, however, is that most published work evaluates sentiment models
in isolation on benchmark datasets, without combining data ingestion, entity
resolution, explainability, and price correlation into a reproducible, user-facing
product. That gap is what this project is designed to address.

\section{Problem Definition}

Despite considerable progress in financial NLP, existing approaches have several
compounding limitations. The key issue is that many systems draw from a single
data source (Twitter or news feeds), producing sentiment signals that reflect
only one segment of the investor population \citep{tetlock2007giving,
bollen2011twitter}. Retail sentiment on Reddit, for example, can diverge
substantially from institutional-oriented news sentiment during the same market
event, as demonstrated by Long et al.\ \citep{long2023reddit} during the
GameStop rally, and ignoring either channel distorts the overall picture.

A second limitation is opacity. Widely used classifiers output a three-class
label (positive, neutral, negative) with a single confidence score, a
representation that discards the behavioural richness of investor discourse.
An investor expressing ``cautious optimism'' differs meaningfully from one
expressing ``confident bullishness'', yet both receive the same
\textit{positive} label. Without knowing \textit{why} a model reached its
classification, users cannot validate outputs, detect biases, or calibrate
trust. \citeauthor{doshivelez2017interpretable}
(\citeyear{doshivelez2017interpretable}) make this requirement explicit, and
\citeauthor{yeo2023xai} (\citeyear{yeo2023xai}) reinforce it in the context
of financial AI. They observe that most XAI work in finance remains confined
to offline research tools rather than deployed user interfaces.

A third gap is the integration of sentiment with price behaviour. While the
literature documents correlations between sentiment and returns
\citep{bollen2011twitter, schumaker2009textmining}, these analyses are
reported in standalone research scripts. No prior work describes an
interactive, ticker-agnostic tool that combines Pearson, Spearman, lag
analysis, Granger causality \citep{granger1969}, rolling correlation, and
market-adjusted metrics within a single user-configurable interface.

To address these gaps, I developed an \textit{Investor Sentiment Dashboard} that
(i) aggregates sentiment across multiple heterogeneous sources, (ii) introduces a
novel finance-specific emotion taxonomy enriching FinBERT outputs with six
behavioural dimensions, (iii) provides an ABSA module for traceable justification,
(iv) surfaces LIME explanations within an interactive React dashboard, and (v)
integrates a sentiment-price correlation engine with Granger causality testing.

\paragraph{Scope note.}
This is a research artefact rather than a regulated financial product. I designed the React interface with production-quality engineering in mind (clear navigation, responsive charts, and explicit Methodology and Legal pages). This allows the system to credibly serve as an investor analytics prototype or the basis for a future pilot study. Where this report refers to \emph{investor-facing} or \emph{production-style} qualities, it means product-minded design that could support user trials or a journal demonstration. It is not a claim of FCA authorisation, suitability advice, or commercial SLAs.

\section{Aims and Objectives}
\label{sec:aims_objectives}

My overarching aim was to design, implement, and evaluate a transparent,
reproducible, end-to-end system for analysing and visualising investor sentiment
across heterogeneous online sources.

The specific objectives were to:

\begin{enumerate}
    \item Aggregate sentiment data from Reddit, X (Twitter), and financial news
    via documented official interfaces (PRAW, the X API v2, and NewsAPI), with
    rate-limit management, a historical CSV fallback for reproducible
    evaluation, and automated scheduling.
    \item Apply FinBERT to classify collected text and formally evaluate it against
    a keyword-based baseline on a purpose-built 200-sample benchmark.
    \item Develop a \textit{novel finance-emotion taxonomy} layering six
    domain-specific emotions (fear, optimism, uncertainty, confidence, scepticism,
    and mixed) over FinBERT outputs using Shannon entropy, domain lexicons, and
    financial aspect priors.
    \item Implement a lightweight \textit{Aspect-Based Sentiment Analysis} module
    that identifies and tags evidence sentences across six financial aspect
    categories.
    \item Integrate LIME as a REST endpoint and embed it within the React SPA,
    making model explanations directly accessible to end users.
    \item Build a stock entity resolution pipeline combining spaCy Named Entity
    Recognition with fuzzy string matching to link free-text mentions to canonical
    ticker symbols.
    \item Implement a \textit{sentiment-price correlation engine} supporting
    Pearson correlation, Spearman correlation, configurable lag analysis, Granger
    causality testing, and rolling correlation.
    \item Develop a React SPA built with Vite and Recharts, and deploy the backend
    in both full (ML) and lean (DB-only) modes.
\end{enumerate}

\section{Novel Contributions}
\label{sec:novel_contributions}

The combination of multi-source ingestion, a finance-specific emotion
taxonomy, lightweight ABSA, end-to-end XAI, and an interactive correlation
engine within a single reproducible tool is not present in any prior system I
have been able to identify in the reviewed literature. This project makes
three primary novel contributions:

\begin{enumerate}

    \item \textbf{Finance-specific emotion taxonomy.}
    Prior work on financial sentiment is largely limited to three-class
    (positive/neutral/negative) classifiers \citep{araci2019finbert,
    malo2014financialphrasebank}. Bollen et al.\ \citep{bollen2011twitter}
    proposed multi-dimensional mood indices for general social media, but no
    equivalent has been defined for the vocabulary and behavioural patterns
    specific to investor discourse.
    I designed and implemented a six-class taxonomy (fear, optimism, uncertainty,
    confidence, scepticism, mixed). It is computed from FinBERT's softmax
    probabilities, normalised Shannon entropy, domain lexicons, and financial
    aspect priors in a single inference pass. I was unable to locate a prior
    system that fuses these four signals within one architecture-agnostic
    scoring function. The module is deployed in production, annotating the
    entire 150{,}000-record corpus.

    \item \textbf{Lightweight ABSA for financial enrichment.}
    Full ABSA systems require fine-tuned sequence-labelling models and
    purpose-built annotated datasets \citep{pontiki2014semeval}, making them
    impractical for high-throughput, continuously ingesting pipelines.
    I implemented a keyword-triggered, sentence-level aspect extraction module
    covering six financial aspect categories (revenue, growth, risk, valuation,
    product/AI, and macro/policy). The module operates without any labelled
    training data and attaches evidence sentences to each stored record,
    providing traceable justification for the headline sentiment label.

    \item \textbf{Production-grade XAI integration.}
    \citeauthor{yeo2023xai} (\citeyear{yeo2023xai}) and \citet{biswas2024xai}
    both find that XAI applications in finance typically remain confined to
    offline notebooks, and this pattern holds for the financial sentiment
    systems I reviewed. I embedded LIME \citep{ribeiro2016lime} as a live FastAPI REST
    endpoint and resolved the async/CPU-blocking challenge by dispatching the
    perturbation sampler onto the default thread-pool executor via
    \texttt{asyncio.to\_thread} (the Python~3.9+ high-level wrapper over
    \texttt{loop.run\_in\_executor}). Token-level heatmaps are rendered directly
    within the React frontend, making model explanations accessible to end
    users without any code interaction. LIME diagnostics also produced a
    concrete finding: contextually neutral risk vocabulary systematically
    triggers negative misclassification, providing an actionable target for
    future fine-tuning.

\end{enumerate}

\section{Report Structure}

The remainder of this report is organised as follows:

\begin{itemize}
    \item \textbf{Chapter 2, Literature Review:} reviews prior research on
    financial sentiment analysis, NLP models, explainable AI, and
    sentiment-price relationships, identifying the gaps that motivated this
    project.
    \item \textbf{Chapter 3, Technical Background:} derives the mathematical
    foundations of the methods used: the transformer self-attention mechanism,
    BERT and FinBERT pre-training, the LIME and SHAP explainability frameworks,
    and the Granger causality F-test.
    \item \textbf{Chapter 4, System Design and Methodology:} details the
    five-tier architecture, development methodology, ethical considerations, and
    key design decisions.
    \item \textbf{Chapter 5, Implementation:} describes the technical
    realisation of each system component, with code listings illustrating
    key choices.
    \item \textbf{Chapter 6, Evaluation:} presents quantitative benchmark
    results, LIME qualitative analysis, and system performance measurements.
    \item \textbf{Chapter 7, Discussion:} contextualises findings against the
    literature, reflects on novel contributions, and explores implications.
    \item \textbf{Chapter 8, Conclusion:} summarises key outcomes, reflects on
    personal and professional development, and proposes directions for future
    research and publication.
\end{itemize}
